\chapter{Resultados preliminares}
\label{cap:resultados}

\notaguia{Primeira redação, a partir das Partes 1--3 da visualização (cujos
números foram auditados um a um contra os arquivos de saída) e da Parte~4 do
\texttt{guia\_de\_leitura.md}. As figuras deste capítulo são geradas por
\texttt{qualificacao/fig/gerar\_figuras.py}, que as redesenha na largura da
mancha a partir dos mesmos arquivos de dados que os pipelines gravaram ---
nunca a partir dos PNG de tela de \texttt{outputs/}.}

Este capítulo apresenta os resultados obtidos até a qualificação, em duas
camadas. A primeira descreve o fenômeno bruto: quarenta anos de uso da terra
em Goiás, organizados pelos três atos que a periodização estabeleceu, e o
balanço agregado de para onde os hectares foram. A segunda percorre as quatro
frentes de investigação, na ordem em que uma puxou a seguinte: o padrão
existe? qual é o mecanismo? uma região empurra a outra? e por que a conversão
desacelerou onde desacelerou?

Duas convenções de leitura valem para todo o capítulo. Onde a medida depende
da classe ``Agricultura'' em janela que alcance 2020--2024, o número é
reportado como intervalo entre as duas réguas da decisão D26, e conferido
contra uma fonte que não passa pelo classificador. E onde um deslocamento
tem intervalo de confiança que inclui zero, ele não é reportado em
quilômetros (decisão D19).

\section{Quarenta anos em três atos}
\label{sec:res-atos}

A régua temporal da Seção~\ref{sec:met-periodizacao} devolveu duas fronteiras,
e a Tabela~\ref{tab:quebras} mostra a convergência dos métodos que as
sustentam. A de 2001 é apontada pelo método primário com folga e corroborada
pelos dois métodos de sensibilidade. A de 2020, pelo primário e,
decisivamente, por uma fonte que não passa por classificador algum: a área de
soja plantada que o IBGE recolhe em campo quebra em 2020 sozinha.

\begin{table}[htb]
\centering
\caption[Convergência dos métodos de periodização]{Convergência dos métodos de periodização sobre as duas fronteiras.}
\label{tab:quebras}
\small
\begin{tabular}{lll}
\toprule
\textbf{Método} & \textbf{$\sim$2001} & \textbf{$\sim$2020} \\
\midrule
sup-F multivariado (primário) & $F = 62{,}2$ & $F = 21{,}5$ \\
Rodionov/STARS (sensibilidade) & 2004--2006 & --- \\
Divergência entre matrizes de transição & pico em 2003 & pico em 2019--2022 \\
Soja plantada, IBGE (âncora independente) & --- & quebra em 2020 \\
\bottomrule
\end{tabular}
\fontefig{elaboração própria a partir do MapBiomas Coleção 10.1 (Pipeline~\#29).}
\end{table}

A caracterização global dos atos vem de uma análise de intensidade, com
teste de Kruskal-Wallis: os três atos diferem em taxa de mudança
($H = 20{,}3$; $p < 0{,}001$), num teste único que compara os três períodos
simultaneamente, e não em detecções separadas por fronteira. Esse número
\emph{descreve} o contraste; ele não o valida. As fronteiras foram estimadas
dos mesmos dados, maximizando precisamente a diferença que o teste depois
mede, e o teste supõe independência entre observações que são taxas anuais
autocorrelacionadas. A validação que vale é externa, e vem adiante: a série
de soja da PAM quebra em 2020 por conta própria, sem nada saber desta
partição.

Uma candidata a quarto período, em torno de 2005--2006, foi examinada e não
adotada. Aqui também convém expor a instabilidade em vez de escondê-la: a
significância da quebra depende do ano de corte testado --- $p = 0{,}10$ em
2004, $p = 0{,}046$ em 2005, $p = 0{,}12$ em 2006 ---, e é justamente por a
fronteira não se firmar em nenhum ano específico que ela não foi adotada, e
não porque um único teste tenha dado $p = 0{,}060$.

Uma ressalva acompanha a fronteira de 2020. Recontada a matriz de transições
com o ``Mosaico de Usos'' como grupo próprio --- consequência da decisão
descrita na Seção~\ref{sec:met-decisoes} ---, o pico do indicador de
divergência migrou de 2020 para 2022, e os quatro métodos não apontam mais
exatamente o mesmo ano. A fronteira se sustenta porque a âncora que a define
é a série de campo, imune ao classificador.

\subsection{Ato I: a pastagem como herança (1985--2000)}

O primeiro mapa da série já mostra uma fronteira em movimento, e não um
ponto de partida natural. Em 1985, quase um terço de Goiás já é pastagem ---
legado da ocupação dirigida do Cerrado que começa na década anterior e que a
série não alcança. A instabilidade macroeconômica do período desestimula o
investimento agrícola de prazo longo, e a lavoura permanece onde já estava:
manchas pontuais no sudoeste. O que domina o centro-sul é pasto; o que
resiste no norte e no nordeste, sobretudo no Vão do Paranã, é Cerrado.

\subsection{Ato II: expansão e intensificação (2001--2019)}

O ato começa com a moeda estabilizada e a exportação de grãos já
desonerada. O que é novo vem de fora: a entrada da China na Organização
Mundial do Comércio no fim de 2001, a sistematização do crédito rural a
partir de 2002--2003 e o super-ciclo de preços de commodities. A agricultura
goiana passa de 9,3\% para 16,0\% do território ao longo do período. No
mapa, o sudoeste --- Rio Verde, Jataí, Mineiros --- se converte à lavoura
sobre o pasto. A pastagem atinge seu pico por volta de 2003, perto de
14,8~Mha, e passa a ceder área. A pecuária segue dominante em extensão, mas
a dinâmica econômica passa à soja.

\subsection{Ato III: conversão acelerada sob rótulo ambíguo (2020--2024)}

Este é o ato em que a medida engana. Lida apenas pela classe ``Agricultura''
do satélite, a conversão parece ter parado: no Sul goiano, a taxa anual do
fluxo de pastagem para lavoura cai 88\% em relação ao Ato~II. A área de soja
plantada que o IBGE recolhe em campo cresce 38\% no estado na mesma janela.
Recontada pela régua que soma as duas classes --- a pergunta mais grossa,
robusta à reetiquetagem por construção ---, essa mesma taxa no Sul não cai:
acelera 51\%. O que mudou foi o rótulo, não o território.

O que o pasto perde, a lavoura ganha: o recuo da pastagem quase quadruplica
de velocidade, de 0,07 para 0,27~Mha por ano. A vegetação natural, por sua vez,
não muda de ritmo --- perde, em saldo líquido, cerca de 0,06~Mha por ano
tanto no Ato~II quanto no Ato~III ---, mas passa a ceder ao norte, onde ainda
há Cerrado em pé.
Goiás chega a 2024 com 34,9\% do território em vegetação natural, sem que a
série dê sinal de estabilização até aqui. O que mudou não foi a velocidade
da perda, e sim o seu endereço. A Figura~\ref{fig:cobertura} mostra esse
rearranjo no mapa, nos cinco cortes temporais.

\begin{figure}[htb]
\centering
\caption[Cobertura e uso da terra em cinco cortes, 1985--2024]{Cobertura e uso da terra em Goiás nos cortes de 1985, 1995, 2005,
2015 e 2024, em sete classes.}
\label{fig:cobertura}
\includegraphics[width=\textwidth]{cap4_painel_cobertura.png}
\fontefig{elaboração própria a partir do MapBiomas Coleção 10.1
(Pipeline~\#10). Os cinco recortes compartilham legenda, projeção, escala e
orientação; por isso a régua e a rosa-dos-ventos aparecem uma única vez, e
valem para todos. Cada pixel mede 387~m. O vazio a leste do centro é o
Distrito Federal, que não integra Goiás e por isso não é classificado.}
\end{figure}

\subsection{O balanço de quarenta anos}

A Tabela~\ref{tab:balanco} resume a transformação agregada, e o saldo líquido
é exatamente o que esconde o caminho dos hectares: duas classes --- pastagem
e Mosaico --- parecem estáveis no saldo e sobem e descem no percurso.

\begin{table}[htb]
\centering
\caption[Balanço agregado de uso e cobertura, 1985--2024]{Balanço agregado de uso e cobertura da terra em Goiás, 1985--2024.}
\label{tab:balanco}
\small
\begin{tabular}{p{3.4cm}p{3.0cm}p{7.2cm}}
\toprule
\textbf{Classe} & \textbf{1985 $\to$ 2024} & \textbf{Leitura} \\
\midrule
Vegetação natural\footnotemark[1] & 17,65 $\to$ 11,88 Mha ($-5{,}8$) &
De 51,9\% a 34,9\% do território \\
\addlinespace
Agricultura & 1,17 $\to$ 5,58 Mha ($\times 4{,}8$) &
A soja cresce $\times 12$ em área (0,37 $\to$ 4,50 Mha) e $\times 13$ em
produção (1,3 $\to$ 17,0 Mt) \\
\addlinespace
Pastagem & 10,98 $\to$ 11,99 Mha ($+1{,}0$) &
Saldo enganoso: a trajetória é em U invertido, com pico próximo de 14,8 Mha
em 2003 \\
\addlinespace
Lotação bovina & 1,45 $\to$ 1,94 cab/ha &
O rebanho cresce 46\% (15,9 $\to$ 23,2 milhões de cabeças) e a área de pasto
apenas 9\%: a pecuária se intensifica \\
\addlinespace
Mosaico de Usos & 3,63 $\to$ 3,59 Mha ($\approx 0$) &
O saldo é nulo, o caminho não: cai de 10,7\% do estado a 6,1\% em 2019 e
volta a 10,5\% em 2024 \\
\bottomrule
\end{tabular}
\fontefig{elaboração própria a partir do painel municipal (Pipeline~\#25),
MapBiomas Coleção 10.1.\\
\textsuperscript{1}~Vegetação natural, nesta tabela, soma formação florestal,
formação savânica, campo nativo \textbf{e campo alagado}. A
Figura~\ref{fig:sankey} adota o agrupamento em sete classes do censo de
pixels, que classifica o campo alagado em ``Outros''; daí a vegetação
natural aparecer lá como 16,96~Mha em 1985, e não 17,65~Mha --- a diferença
é o campo alagado, 0,70~Mha (a subtração dos valores já arredondados dá
0,69). As duas contas são a mesma medida sob duas convenções de classe, e
não uma divergência de dado.}
\end{table}

O cruzamento pixel a pixel entre 1985 e 2024 mostra de onde veio a terra que
cada classe ganhou. Três fluxos desenham a marcha: a conversão de vegetação
natural em pastagem, com 4,10~Mha, é o maior da série --- desmatamento
histórico para pecuária extensiva, concentrado sobretudo nos anos 1980 e
1990; a conversão de pastagem em agricultura, com 2,72~Mha, é a maior
transição entre classes já manejadas, e é ela que permite a expansão da
lavoura sem desmatamento direto; e a conversão direta de vegetação natural
em agricultura, com 1,29~Mha, menor que a rota via pasto, mas existente e
com peso crescente no nordeste após 2010. A pastagem é, portanto, o elo
intermediário da maior parte da entrada em agricultura --- 2,72 dos
4,01~Mha, ou 68\%.

Um quarto fluxo --- 1,00~Mha de vegetação natural para Mosaico de Usos ---
foi deliberadamente mantido fora da conta de conversão antrópica, e a razão
é externa ao próprio dado: incluí-lo faria a razão entre a medida deste
trabalho e a do PRODES/INPE saltar de 1,00 para 1,35, isto é, um terço a
mais de desmatamento do que o sistema oficial registra. Como a classe tem
transição de ida e de volta, tratá-la como perda definitiva superestima.

\begin{figure}[htb]
\centering
\caption[Cruzamento pixel a pixel do uso da terra, 1985--2024]{Cruzamento pixel a pixel do uso da terra entre 1985 e 2024, em sete
classes, com larguras proporcionais à área.}
\label{fig:sankey}
\includegraphics[width=\textwidth]{cap4_sankey.pdf}
\fontefig{elaboração própria a partir da recontagem sobre o cubo do censo de
pixels (Pipeline~\#12B), MapBiomas Coleção 10.1. Fitas inferiores a
0,05~Mha foram omitidas para leitura; os quatro fluxos discutidos no texto
estão rotulados. O agrupamento em sete classes classifica o campo alagado em
``Outros'', e é só por isso que a vegetação natural aparece aqui com
16,96~Mha em 1985 contra os 17,65~Mha da Tabela~\ref{tab:balanco}, que o
conta como natural --- ver a nota daquela tabela.}
\end{figure}

\section{O padrão: a marcha ao norte}
\label{sec:res-perna1}

A primeira frente responde à pergunta mais básica: a fronteira agropecuária
de Goiás se moveu, e para onde? A medida é o centro de massa de cada uso da
terra --- o ponto de equilíbrio geográfico ponderado pela quantidade da
variável ---, acompanhado ano a ano sobre as 166 AMCs. O resultado é
uniforme: tudo marchou ao norte (Tabela~\ref{tab:centros}).

\begin{table}[htb]
\centering
\caption[Deslocamento do centro de massa ao norte, 1985--2024]{Deslocamento do centro de massa ao norte, 1985--2024, com
intervalo de confiança de 95\% por \emph{bootstrap} das AMCs
($B = 2000$).}
\label{tab:centros}
\small
\begin{tabular}{lrll}
\toprule
\textbf{Variável} & \textbf{$\Delta$ Norte} & \textbf{IC 95\%} & \textbf{Veredito} \\
\midrule
Pastagem          & $+77{,}6$ km & $[+54{,}7;\ +98{,}2]$ & robusto \\
Rebanho bovino    & $+66{,}9$ km & $[+47{,}2;\ +84{,}5]$ & robusto \\
Agricultura       & $+65{,}2$ km & $[+43{,}5;\ +94{,}6]$ & robusto \\
Vegetação natural & $+7{,}6$ km  & $[-0{,}5;\ +15{,}6]$  & \textbf{ancorada} \\
\bottomrule
\end{tabular}
\fontefig{elaboração própria (Pipeline~\#32) sobre as 166 AMCs; intervalos de confiança de 95\% por \emph{bootstrap} das AMCs ($B = 2000$).}
\end{table}

Três leituras decorrem da tabela. O deslocamento da pastagem, do rebanho e
da agricultura é grande e está longe de zero. A vegetação natural não
acompanhou: seu intervalo inclui zero, e por isso ela é descrita como
ancorada, e não como tendo andado oito quilômetros --- o que seria ler ruído
como sinal. E existe um gradiente latitudinal persistente: em todos os
quarenta anos, a lavoura permanece cerca de 120 a 130~km ao sul do pasto e do
rebanho --- a expressão métrica dos ``dois Goiáses''.

O tamanho exato desse vão depende da régua: em 2024 ele é de 135~km pela
classe do satélite, 111~km pela régua da união e 122~km pela soja que o IBGE
mede em campo, grande e persistente em qualquer uma das três. A
Figura~\ref{fig:centromassa} traz as trajetórias completas, com a faixa de
confiança de cada série.

\begin{figure}[htb]
\centering
\caption[Centro de massa de quatro camadas]{Centro de massa de quatro camadas: posição em 1985 e 2024 sobre a
malha das AMCs (a) e latitude ano a ano com a faixa de \emph{bootstrap} (b).}
\label{fig:centromassa}
\includegraphics[width=\textwidth]{cap4_centro_massa.pdf}
\fontefig{elaboração própria (Pipeline~\#32) sobre as 166 AMCs; faixa de
confiança de 95\% por \emph{bootstrap} das AMCs ($B = 2000$). Mapa em
projeção Policônica do Brasil (SIRGAS 2000, EPSG:5880), a mesma em que os
centroides são calculados. O mapa mostra o
estado inteiro de propósito: a marcha líquida é de cerca de 80~km num estado
com cerca de 790~km de extensão norte--sul, e um zoom na nuvem de
trajetórias faria um sinal
pequeno, porém robusto, parecer maior do que é. A série da agricultura
aparece \textbf{pontilhada a partir de 2019} porque, dali em diante, parte da
lavoura recém-convertida passa a receber o rótulo ``Mosaico de Usos'' e deixa
de entrar neste centroide (decisões D25 e D26); o achatamento visível no
Ato~III é, em parte não medida, artefato de rótulo.}
\end{figure}

\subsection{Sensibilidade da marcha à régua de medida}

O centroide da agricultura pesa cada região pela área que o satélite rotula
como agricultura --- e, a partir de 2021, parte da conversão recente passa a
receber o rótulo de Mosaico: exatamente a lavoura nova, a da fronteira, deixa
de entrar na conta. A linha achata depois de 2019, e a leitura tentadora
seria que a agricultura parou de subir.

O teste é possível porque nem toda medida do painel passa pelo
classificador. Entre 2019 e 2024, quatro medidas indicam deslocamento ao
norte --- três delas entre 10 e 13~km, a quarta 4,4~km --- e a quinta indica
que nada se moveu. A quinta é precisamente a única cujo rótulo mudou no
período. O rebanho e a soja plantada, medidos em campo, marcham.

Uma segunda evidência fecha o argumento: localizada no mapa, a área que mudou
de rótulo não está espalhada --- o centro dessa massa fica 46~km ao norte
do centro da agricultura visível, e seu crescimento por região acompanha o
da soja do IBGE quase ponto a ponto (correlação de 0,84; 1,52 contra
1,53~milhão de hectares). A conversão que o satélite deixou de nomear está
acontecendo na fronteira. O erro de medida, portanto, aponta contra a tese e
não a favor dela: a régua crua faz a marcha parecer mais fraca do que foi.

Cabe uma distinção que a Seção~\ref{sec:res-perna3} tornará central: dizer
que a lavoura nova está ao norte não é dizer que a lavoura do Sul empurrou
alguém para o norte. A lavoura está subindo junto, não empurrando de trás. O
deslocamento da fronteira agrícola é medido aqui; o de uma região pela outra
foi testado, e não foi encontrado.

A régua da união, porém, não vale para a série longa. Aplicada aos quarenta
anos, ela devolve $-60$~km, isto é, a agricultura andando para o sul --- e o
número não significa o que parece. Em 1985 a classe Mosaico já existia e era
extensa --- 3,63~Mha, 10,7\% do estado ---, situada bem ao norte, e não era
lavoura sob outro nome: era terra de uso misto indiferenciado. Ao longo da
série ela encolhe enquanto a lavoura cresce ao sul, de modo que somar as duas
classes nas duas pontas é comparar objetos diferentes: mede-se sobretudo o
Mosaico antigo se dissolvendo.

A confirmação é espacial: o Mosaico ocupa praticamente a mesma área nas duas
pontas da série, mas as manchas quase não se sobrepõem --- a correlação
espacial entre as células é de apenas 0,15. A de 1985 se dissolveu; a de 2024
surgiu em outro lugar. Por isso a régua da união entra como teto de uma
janela curta, 2019--2024, onde a migração de rótulo de fato ocorre, e nunca
como régua de série longa.

\subsection{Controles de escala e a camada do fogo}

A marcha sobrevive a quatro verificações independentes. Não é artefato de
recorte: refeito pixel a pixel, sem malha administrativa alguma, o
deslocamento reaparece a um ou dois quilômetros do valor original (pastagem
$+79{,}2$ contra $+77{,}6$~km; agricultura $+66{,}9$ contra $+65{,}2$~km).
Não é artefato de agregação de classes: aberta em formações, a vegetação
mostra que a resistência ao norte é da floresta ($+8{,}7$~km, IC
$[+2{,}5; +15{,}1]$), ao passo que o limite do campo nativo se deslocou para
o norte --- direção robusta, magnitude imprecisa, com intervalo largo --- e a
formação savânica tem intervalo que inclui zero. Foi uma autocorreção: o
número médio do agregado escondia comportamentos distintos por dentro. Toda a
marcha vem com barra de erro, pela decisão D19.

E não é deslocamento genérico, o que dois controles estabelecem. A área
urbana não acompanha: ela anda 8~km para o sul (IC $[-19; -1{,}5]$), no
sentido oposto ao da fronteira. E o leite --- pecuária da bacia consolidada
do Sul --- sobe menos da metade do que sobe o rebanho total ($+30$ contra
$+67$~km), de modo que o vão entre os dois centros mais que dobra, de 30~km
em 1985 para 68~km em 2024. Quem marcha é a pecuária de corte na fronteira,
e não o mapa inteiro subindo.

Há ainda uma quinta camada, e ela vai à frente das demais: o centro de massa
do fogo em vegetação natural fica ao norte do centro da conversão em todos
os 39 anos em que as duas séries coexistem --- a de fogo começa um ano
depois da de cobertura ---, com dianteira média de 73~km, sem nunca trocar
de lado. É uma
vanguarda geográfica, e apenas isso: a leitura de que o fogo abriria caminho
foi testada e não se sustenta no tempo --- a liderança temporal desaparece
sob transformação logarítmica, inverte quando se removem os anos de seca de
1985 e 2010, e o teste de precedência agregado é nulo. O fogo marca onde a
fronteira está, não quando ela vai andar.

\section{O mecanismo: duas lógicas simultâneas}
\label{sec:res-perna2}

Se a fronteira inteira marchou, que conversão a moveu, e em que parte do
estado? A marcha é o saldo de duas conversões distintas, e elas não estão no
mesmo lugar. No Sul, a transição dominante é de pastagem para lavoura:
intensificação, terra que troca de função. No Norte, é de vegetação natural
para pastagem: fronteira, terra nova sendo aberta.

A medida que sustenta essa divisão é justamente a que não depende da classe
ambígua, porque nem sua origem nem seu destino passam por ela: a conversão
de vegetação natural em pastagem cai 49\% no Sul Goiano entre os Atos~II e
III e apenas 13\% no Norte, na malha das cinco mesorregiões.

\subsection{A idade da pastagem como medida de intenção}

Saber que o Sul converte pasto em lavoura não diz se aquele pasto era etapa
planejada de um sistema agrícola ou terra parada há duas décadas que enfim se
tornou lucrativa --- duas economias diferentes com o mesmo nome no mapa. A
idade da pastagem no instante da conversão é a aproximação mais próxima da
intenção do produtor que a série permite observar.

O censo de pixels descrito na Seção~\ref{sec:met-decisoes} contabiliza
44,6~milhões de conversões de pastagem para lavoura em Goiás --- 3,8~Mha,
11,2\% do estado. Destas, 16,0~milhões têm idade conhecida; as demais já
eram pastagem em 1985 e têm idade truncada pelo início da série.

A distribuição não é composta por um único tipo de pastagem. O ajuste de
mistura sobre o censo completo separa duas populações: uma jovem e
concentrada ($\mu \approx 4$ anos, $\sigma \approx 1{,}6$ ano, 33\% da
massa) e outra velha e larga ($\mu \approx 16$ anos,
$\sigma \approx 7{,}5$ anos, 67\% da massa), com cauda que alcança os
quarenta anos da série. A forma não tem dois picos visíveis: tem um pico
jovem e um ombro longo --- sendo cerca de cinco vezes mais dispersa que a
primeira, a segunda população produz platô, e não segunda espiga. O que a
estabelece não é a inspeção visual, e sim o fato de uma população única não
dar conta da forma: ajustada sozinha, ela devolve $\mu \approx 12$ anos e
erra as duas pontas, subestimando o pico jovem e a cauda velha.

A leitura de cada grupo apoia-se no que havia no pixel antes de ele virar
pastagem --- e essa leitura não depende de ajuste algum. Separados apenas
pela origem, os eventos se ordenam sozinhos: o pasto que veio de lavoura
volta a ser lavoura com mediana de cinco anos, e o que veio de vegetação
natural dura o dobro, mediana de treze anos. Nas conversões de pasto com até
três anos, 45\% já tinham sido lavoura antes --- lavoura, pasto, lavoura: é
rotação. Nas de pasto com trinta anos ou mais, a lavoura anterior
praticamente desaparece, com 1\%: esse pasto veio de vegetação natural ou de
Mosaico de Usos, em proporções semelhantes --- Cerrado, pasto, lavoura: é
reserva ativada.

Cabe aqui uma ressalva, porque o Mosaico é justamente a classe em que o
classificador não separou usos: metade da origem do pasto velho fica, por
construção, indeterminada entre vegetação natural e uso misto anterior. O
que o dado estabelece com firmeza é o \emph{contraste} entre as duas pontas
--- 45\% contra 1\% de lavoura anterior ---, e não a atribuição de cada
pixel. Na mesma linha, a primeira população é compatível com sistemas de
integração lavoura-pecuária, mas não os comprova: o satélite vê o capim, não
vê se havia gado nele. A Figura~\ref{fig:idade} traz a distribuição de
idades e a decomposição por origem anterior.

\begin{figure}[htb]
\centering
\caption[Idade da pastagem no momento da conversão para lavoura]{Idade da pastagem no momento da conversão para lavoura: a forma da
distribuição contra os ajustes de uma e de duas componentes (a) e a
composição da origem anterior do pixel por idade (b).}
\label{fig:idade}
\includegraphics[width=\textwidth]{cap4_idade_pastagem.pdf}
\fontefig{elaboração própria a partir do censo de pixels (Pipeline~\#28C),
MapBiomas Coleção 10.1. Entram apenas os 16,0~milhões de eventos cuja idade é
conhecida; os demais já eram pastagem em 1985 e têm idade truncada pelo
início da série. O painel~(a) mostra por que a segunda população não se lê
por inspeção visual: sendo cerca de cinco vezes mais larga que a primeira,
ela produz um ombro, e não um segundo pico. O painel~(b) recupera a mesma
divisão sem ajustar modelo algum.}
\end{figure}

\subsection{Decomposição regional da idade}

A curva soma quarenta anos de conversões em todo o estado: talvez não haja
duas populações convivendo em lugar algum, e sim regiões ou épocas diferentes
somadas. As duas objeções foram testadas isolando cada recorte.

\textbf{Não são regiões somadas.} Percorridas as cinco mesorregiões, o
desenho é praticamente o mesmo. As cinco são bimodais, e também o são todas
as 166 AMCs na régua imune à mudança de rótulo (na régua estrita, 162 de
164, e as duas exceções falham pelo mesmo lado: peso do grupo jovem baixo
demais para o teste, não ausência de dois grupos). A mesorregião explica
0,5\% da variação da idade.

\textbf{Não são épocas somadas.} Isoladas por período, as duas populações
permanecem: no Ato~II sozinho --- 32,5~milhões de conversões pela régua da
união, o maior bloco da série --- e no Ato~III sozinho, nas cinco
mesorregiões de cada um, em 10 de 10 células de período por região (9 de 10
pela régua estrita). Convém não somar esses 32,5 milhões aos 16,0 milhões
citados acima: aqueles são eventos com idade conhecida na régua estrita,
estes são eventos na régua da união, e as duas contagens medem populações
diferentes. O Ato~I fica fora do teste, e a
razão é instrutiva: ele é unimodal em toda parte e não poderia ser outra
coisa. Uma conversão em 1995 ocorre sobre um pasto que tem, no máximo, dez
anos, porque a série começa em 1985; para observar a conversão de um pasto
de dezesseis anos --- o modo da população velha --- é preciso chegar a 2001.
No primeiro período, a população velha é invisível, não ausente.

Aquele 0,5\% é o que restou da autocorreção mais extensa do trabalho,
apresentada aqui como resultado superado --- a leitura que seria publicada se
um teste específico não tivesse sido feito. A afirmação superada era que o
Sul converteria pasto jovem (mediana de nove anos) e o Norte, pasto velho
(dezesseis anos). O defeito não estava na conta, e sim no que entra nela: a
idade só é medida quando o destino do pixel recebe o rótulo de agricultura e,
no fim da série, parte da conversão passa a ser rotulada como Mosaico. Esses
eventos saem da amostra sem aviso, e não saem por igual no território: saem
mais ao norte. O gradiente era, em boa parte, essa seleção.

Refeita a análise com a pergunta grossa --- imune à reetiquetagem por
construção ---, a amplitude Sul--Norte cai de sete anos para dois, a ordem
das regiões se embaralha e a parcela de variação atribuível à mesorregião
despenca de 3,7\% para 0,5\%. O mesmo ocorre com a forma da distribuição:
na régua estrita, Norte e Noroeste exibem um vale nítido que Sul, Centro e
Leste não têm --- o vale do Noroeste tem profundidade de 42\% ---, e as
cinco regiões se separam em dois blocos. Fechado o viés, a distância entre a
forma do Sul e a do Norte cai de 0,22 para 0,02, e o vale do Noroeste, de
42\% para 6\%.

O que não mudou sob a régua nova foi a coexistência das duas populações. O
achado que sobrevive é, portanto, mais forte do que o que caiu: quanto menos
a região explica, mais universal se mostra a convivência dos dois
mecanismos.

\subsection{Sistema produtivo e crédito como candidatos a explicação}

Descartados os dois eixos óbvios, restavam duas famílias de explicação, cada
uma com uma previsão verificável. A primeira é estrutural: a idade seria
decidida pelo sistema produtivo instalado. O Censo Agropecuário não tem
variável de integração lavoura-pecuária, e o indicador mais próximo
disponível é o plantio direto --- que é prática de conservação de solo, não
de integração com pecuária, ainda que ambas costumem andar juntas na
lavoura tecnificada do Cerrado. A segunda é de fluxo: a idade seria decidida
pelo choque do momento, com o crédito ativando reserva antiga.

Nenhuma das duas explica a idade da pastagem. Para o plantio direto, a
correlação simples confirma a previsão, mas o indicador também se concentra
no Sul; segurando o gradiente espacial (decisão D14), a associação encolhe
cerca de um terço, fica na fronteira da significância sem cruzá-la e não
sobrevive à correção para testes múltiplos. O veredito é
\emph{não estabelecido}, e a distinção importa: um nulo limpo encerraria a
pergunta, um indeterminado não encerra.

Para o crédito rural, a associação é a única que resiste ao controle
espacial, à correção para testes múltiplos e à troca de régua --- na régua
imune ela se fortalece, de $+0{,}22$ para $+0{,}30$. Mas aponta na direção
contrária à esperada: onde o crédito cresceu mais, o pasto convertido é mais
velho, o que é compatível com reserva sendo ativada por dinheiro novo e não
com crédito financiando rotação. Duas ressalvas grandes acompanham o
resultado: ele só existe quando os anos recentes entram --- no recorte até
2019 é praticamente nulo --- e o mecanismo não foi investigado.

A escolha entre girar o capim em três anos ou deixá-lo trinta acontece,
portanto, abaixo da escala que estes dados alcançam: ela não se lê na região,
não se lê no sistema produtivo do município e mal se lê no crédito que chegou
até ele.

\section{O teste do deslocamento e o \emph{drive} comum}
\label{sec:res-perna3}

As duas frentes anteriores entregam dois fatos que, postos lado a lado,
quase escrevem sozinhos uma conclusão: a fronteira inteira subiu ao norte, e
subiu fazendo coisas opostas em cada metade do estado. Somados, sugerem que
o pasto expulso do Sul foi reaparecer no Norte, sobre o Cerrado --- um
deslocamento indireto intra-estadual, com a soja do Sudoeste na origem e a
vegetação do Norte na conta. É a hipótese que mais favorecia a tese inicial
deste trabalho.

Ela tem um problema de lógica antes de ter um problema de dados. Uma
coincidência no espaço não é um mecanismo: dois processos independentes
movidos pela mesma força externa desenham exatamente o mesmo mapa que um
empurrão desenharia. Separar as duas leituras exige procurar uma assinatura
que só o empurrão produz.

\subsection{Efeito de vizinhança e efeito local}

São duas as assinaturas, e ambas são específicas o bastante para poder
falhar --- o que as torna teste, e não ilustração. No tempo, se a lavoura do
Sul empurra, ela precisa mover-se primeiro: a expansão no Sul deve anteceder
o avanço do pasto no Norte, e não apenas acompanhá-lo. No espaço, um empurrão
tem direção: se a lavoura dos vizinhos ao sul cresce e isso expulsa pasto
para cima, o pasto local deve crescer --- um coeficiente $\theta$ que a
hipótese exige seja positivo, com os vizinhos ao norte entrando como placebo.

O teste espacial foi executado em três réguas de medida, cruzadas com duas
janelas e dois desfechos --- doze especificações. \textbf{As doze estimativas
são negativas.} Nenhuma cai na região que a hipótese exige; todas vão na
direção oposta. Quando a lavoura dos vizinhos ao sul cresce, o pasto local
diminui, em vez de aumentar. Não se trata de um p-valor rejeitando uma tese:
trata-se da ausência de um sinal que a hipótese exigia, procurado em doze
recortes e não encontrado em nenhum.

O teste temporal foi executado em vinte e quatro combinações de régua,
janela, desfecho e defasagem. O menor p-valor é 0,078; nenhuma combinação é
significativa. Este nulo tem limites que devem ser declarados: com 38
diferenças anuais, o teste descarta um empurrão forte (poder aproximado de
93\%), mas não um moderado (poder aproximado de 48\%). Sozinho, não
fecharia a questão; somado às doze especificações espaciais, completa o
quadro.

Rodado no sentido inverso, esse mesmo teste temporal produzia um resultado
forte: o Norte anteciparia o Sul, com $p = 0{,}0007$. Se real, ele não
salvaria a hipótese do empurrão --- inverteria a tese inteira. Investigado,
era regressão espúria, e o diagnóstico veio em três passos, nenhum deles ``a
amostra é pequena''. Primeiro, a régua estava errada para o material: o teste
comparava uma série já estabilizada com outra que permanece instável mesmo
depois de convertida em variações anuais --- a montagem clássica que fabrica
precedência. Segundo, o método adequado a séries integradas apaga o resultado
nas duas direções: $p = 0{,}45$ no sentido Norte--Sul e $p = 0{,}25$ no
sentido Sul--Norte. Dizer os dois importa, porque isso também retira qualquer
autorização para afirmar que o Sul lidera: o veredito é \emph{sem líder}, não
``o outro lado venceu''. Terceiro, e mais limpo, os placebos acendem onde não
há mecanismo: o pasto do Norte ``antecede'' o pasto do próprio Sul com
exatamente o mesmo p-valor do resultado original, o que nenhuma economia
explicaria. Uma precedência que aparece em qualquer par de séries suaves não
é mecanismo, é defeito de instrumento.

Deste episódio saiu a regra D16, aplicada retroativamente aos demais
pipelines: nenhum resultado de precedência em série agregada é lido como
causal antes de passar por diagnóstico de estacionariedade, método robusto à
integração e bateria de placebos. A Figura~\ref{fig:espacial} reúne as doze
estimativas do termo de vizinhança.

\begin{figure}[htb]
\centering
\caption[As doze especificações do teste espacial]{As doze especificações do teste espacial do empurrão, com o
intervalo de confiança de cada uma e a zona que a hipótese exigiria.}
\label{fig:espacial}
\includegraphics[width=\textwidth]{cap4_teste_espacial.pdf}
\fontefig{elaboração própria (Pipeline~\#49). $\theta$ é o coeficiente do
termo de vizinhança ao sul; as especificações de placebo (vizinhos ao norte)
não estão representadas --- e sobre elas cabe um registro, porque não são
uniformemente nulas: numa das seis células (régua da soja do IBGE, janela
plena) o termo de vizinhança ao norte é significativo a 5\%
($\beta = -0{,}054$; $p = 0{,}032$), justamente onde o termo ao sul
correspondente é nulo. Como o argumento da frente é que um empurrão teria
direção, esse acendimento do placebo direcional enfraquece a premissa e é
mais um motivo para não ler as doze estimativas como um achado.
A faixa sombreada é desenhada porque o resultado é
a \emph{ausência de um sinal que a hipótese exigia} --- e não um p-valor
rejeitando uma tese: as doze estimativas caem do lado oposto ao exigido.}
\end{figure}

O mesmo desenho revela o que ocupa o lugar do empurrão, e o termo precisa ser
preciso: o coeficiente $\theta$ mede o efeito dos vizinhos --- se a lavoura
cresce ao sul, o pasto daqui cresce? --- ao passo que o efeito local mede o
que ocorre dentro do próprio município --- se a lavoura cresce aqui, o pasto
daqui diminui? É a diferença entre empurrão, à distância, e substituição, no
mesmo lugar.

A substituição local é o maior efeito de toda a frente: nas quatro células
medidas em área de satélite o coeficiente vai de $-0{,}52$ a $-1{,}14$, e
$p < 0{,}001$ em cinco das seis células da grade (a exceção é a régua da
soja do IBGE na janela curta, com poucos anos). As duas células restantes,
que regridem sobre a área de soja declarada ao IBGE, ficam em torno de
$-0{,}03$ a $-0{,}07$ --- números que \emph{não} são comparáveis em
magnitude aos primeiros, porque o regressor é outro: um hectare de soja é
subconjunto de um hectare de agricultura. O que se compara entre réguas é o
sinal e a significância, não o tamanho.

O coeficiente dobra quando se amplia a medida --- de $-0{,}52$ na régua
estrita a $-1{,}14$ na que soma agricultura e uso misto. Convém ler esse
salto com cuidado, e não como correção: a régua da união é o teto do
intervalo previsto na decisão D26, não a medida certa. E um coeficiente de
$-1{,}14$ não descreve substituição mais completa --- descreve mais de um
hectare de pasto perdido por hectare de lavoura ganho, o que por definição
excede a substituição e inclui algum fluxo adicional não identificado neste
desenho.

\subsection{O \emph{drive} comum incide sobre um gradiente de aptidão}

Se os dois mecanismos não se comandam, por que marcham juntos? Porque ambos
respondem ao mesmo estímulo externo, e o candidato mais forte é o câmbio
real. O mecanismo é direto: Goiás produz soja e carne para exportação, ambas
cotadas em dólar; quando o real se desvaloriza, a mesma tonelada exportada
rende mais reais ao produtor sem que o preço internacional mude; a
rentabilidade de plantar e de criar sobe, e a procura por terra cresce. O
câmbio não cria fazenda --- ele torna compensador avançar.

Esse impulso, porém, não incide sobre terreno uniforme. O gradiente de
aptidão foi medido por um dado físico que nada sabe sobre quem planta o quê:
o mapa de aptidão agrícola da Embrapa, na escala 1:500.000, com 8.284
polígonos recortados em Goiás. Na escala ordinal da fonte, o gradiente
aparece sozinho --- Sul 4,69, Centro 4,47, Norte 4,17 ---, com correlação de
$-0{,}44$ com a latitude ($-0{,}40$ por Spearman; as três regiões diferem
com $p = 0{,}0002$). Medir o gradiente por um dado físico, e não pelo uso da
terra que se quer explicar, é o que evita a circularidade.

O modelo testa a interação entre os dois, e o desfecho da especificação
confirmatória precisa ser dito com todas as letras: é a \textbf{variação
anual do rebanho bovino}, medida em milhões de cabeças, e não uma medida de
área. A previsão é a de um gradiente invertido em relação ao que o senso
comum sugeriria: sob depreciação cambial o rebanho cresce \emph{mais onde a
aptidão edafoclimática é baixa} --- a fronteira do Norte --- e menos no
núcleo apto do Sul, onde a mesma terra tem uso alternativo mais rentável em
lavoura. O coeficiente esperado da interação é, portanto, negativo.

E aqui entra a ressalva que determina o peso do resultado: sob a inferência
apropriada ao desenho --- permutação do \emph{shifter}, conforme a
Seção~\ref{sec:met-instrumental} ---, o achado é \textbf{corroborante, não
estabelecido}. O coeficiente tem a direção prevista ($\beta = -0{,}033$, com
$R^2$ \emph{within} de 0,001), mas o p-valor honesto fica entre 0,07 e 0,13,
sem cruzar o corte de 5\%. O que sustenta a leitura é a especificidade: cada
teste que deveria sair vazio sai vazio. Os placebos de desfecho são nulos
(área urbana, $p = 0{,}34$; água, $p = 0{,}33$); os placebos no tempo também
($p = 0{,}105$ para a liderança de um ano, $p = 0{,}77$ para a de dois); e o
sinal é estável no \emph{jackknife}, em 38 dos 38 anos.

Três qualificações completam o quadro, e todas as três puxam para baixo o
peso do achado. A primeira é que os desfechos de \emph{área} da mesma grade
são nulos: a interação entre câmbio e aptidão não move a variação de área de
pastagem ($p = 0{,}92$ sem defasagem e $0{,}83$ com um ano), e a interação
entre preço da soja e aptidão não move a variação de área de agricultura
($p = 0{,}77$ e $0{,}60$). Como a narrativa desta seção é sobre procura por
terra, é preciso registrar que o único gradiente com sinal aparece numa
contagem de cabeças, e não na terra. A segunda é de multiplicidade: essa
célula vem de uma grade de 192 testes (quatro exposições, quatro
\emph{drivers}, quatro desfechos, três defasagens), e o próprio exame de
sensibilidade mostrou que ela sobrevive ao controle de falsas descobertas
numa família de 96 testes e não sobrevive numa de 144 --- é, portanto,
sensível ao recorte da família, além de assentar sobre um p agrupado que já
se sabe otimista. A terceira é que o desenho, com efeito fixo de ano, não
poderia mesmo estimar o nível do efeito cambial: ele estima só o gradiente,
e é sobre o gradiente --- não sobre o câmbio como causa --- que a evidência
aqui é corroborante.

Uma precisão sobre esses p-valores, sob pena de se comparar réguas
diferentes: o p do achado é o de permutação, o conservador; os dos placebos
são os agrupados, que a auditoria mostrou serem otimistas. Isso corta a
favor da leitura --- sob a régua conservadora, os placebos ficariam ainda
mais nulos ---, mas não autoriza comparar 0,105 com 0,07--0,13 como se
fossem a mesma medida.

\subsection{Precedência de infraestrutura e crédito}

Uma verificação independente examina se infraestrutura ou crédito puxam a
fronteira. O raciocínio é de simetria: quem puxa costuma estar à frente. Se
silos, bancos ou lavouras puxassem a marcha, seus centros de gravidade
estariam ao norte do pasto e do rebanho; se apenas consolidassem o que já
foi aberto, estariam ao sul.

Nenhuma camada aparece à frente. O centroide do crédito rural fica cerca de
70~km ao sul da pastagem, e esse vão não se fecha --- era 77~km em 2013 e
segue em 69~km em 2024. O da capacidade de armazenagem está ainda mais ao
sul, a cerca de 150~km do pasto, praticamente sobre o núcleo de lavoura do
Sudoeste, a 16~km dele. A cadeia exportadora, medida à parte, co-move sem
liderar. Silo, banco e porto consolidam onde a produção já se instalou. O
resultado é descritivo: diz onde cada camada está, não que ela seja incapaz
de puxar.

\section{A desaceleração e o teto de oferta}
\label{sec:res-perna4}

No Ato~III, o Sul quase deixa de abrir terra nova. A explicação intuitiva
seria o esfriamento do mercado. Ela está errada, e é isso que torna a última
frente interessante: no Ato~III a demanda estava no pico.

Duas precisões mudam a pergunta. A primeira: não foi a lavoura que freou.
Essa leitura seria o artefato de rótulo já desarmado. O que freou é a
abertura de vegetação natural, medida pela conversão de vegetação em
pastagem --- transição cuja origem e cujo destino ficam fora da classe
ambígua. Ela cai 49\% no Sul Goiano e apenas 13\% no Norte, medida na malha
das cinco mesorregiões.

A segunda: não foi o estado que desacelerou, foi o Sul. No agregado de
Goiás, o fluxo \emph{bruto} de conversão de vegetação natural --- a abertura,
antes de descontar o fluxo reverso --- passou de 0,071~Mha por ano no Ato~II
para 0,072 no Ato~III, praticamente o mesmo número. O que mudou foi o
endereço: o fluxo bruto caiu 37\% na região Sul e subiu no Centro e no Norte
--- este na malha de três regiões, e por isso não comparável termo a termo
com os 49\% do parágrafo anterior, que são de outra medida e de outra malha. A marcha não perdeu
velocidade; mudou de lugar. A pergunta desta frente não é, portanto, por que
a fronteira parou --- ela não parou ---, e sim o que a expulsou do Sul.

\subsection{Demanda e oferta na desaceleração do Sul}

Os três sinais de demanda acompanhados desde a frente anterior sobem todos
no Ato~III, e não pouco: o câmbio real vai de 134,5 a 169,0, o preço
recebido pela soja de 104,4 a 186,4 e o crédito rural de Goiás de 14,3 para
24,1 bilhões de reais (série da fonte em reais de 2010, e não na base de
dezembro de 2024 usada no restante do trabalho). A soja efetivamente plantada cresce 38\% na mesma
janela.

O teste que fecha a questão, porém, não é o do estado: é o do próprio Sul,
onde a freada ocorreu. Se o Sul tivesse perdido apetite por terra, teria
parado de converter qualquer terra. Não foi o que houve: no Sul do Ato~III, a
taxa anual de conversão de pastagem em lavoura ou uso misto sobe 51\% em
relação ao Ato~II, e o ritmo de expansão da soja plantada --- os hectares que
o IBGE registra como acréscimo a cada ano --- sobe 244\%, ao mesmo tempo em
que a abertura de Cerrado novo cai pela metade. As duas comparações do
parágrafo anterior e desta, convém frisar, não são a mesma medida: 38\% é o
crescimento do estoque plantado no estado entre as pontas da janela, e 244\%
é a aceleração do acréscimo anual no Sul entre um ato e outro. O Sul não
parou de querer terra. Parou de tirá-la do Cerrado e passou a tirá-la de
pasto já aberto.

A mudança do fluxo entre dois atos separa-se aritmeticamente em duas
parcelas: quanto vem de haver menos Cerrado a converter e quanto vem de
converter a uma taxa diferente o Cerrado que ainda há. A decomposição mostra
que a parcela de estoque é negativa nas três regiões --- em toda parte há
menos Cerrado do que havia, inclusive onde a conversão acelerou --- e que o
que separa as regiões é a parcela residual, que desaba no Sul e sobe no
Centro e no Norte.

O limite dessa figura precisa ser explícito, porque a leitura descuidada dela
diz o contrário do que a evidência sustenta. Apenas 17\% da freada do Sul é,
mecanicamente, ter menos Cerrado; os outros 83\% estão na parcela residual.
E essa parcela \emph{não é demanda}: ela reúne tudo o que não é o tamanho do
estoque --- propensão a converter, atrito de acesso, proteção e a troca de
fonte de terra descrita acima. Se o argumento de oferta dependesse do tamanho
da parcela de estoque, ele seria fraco.

Ele não depende. O argumento se monta por eliminação e por regularidade:
a taxa do Sul caiu enquanto a demanda estava no pico, de modo que o que a
derrubou não é falta de comprador; não é a lei, pelo que se mostra adiante;
a demanda de terra do Sul continuou existindo e foi atendida por pasto já
aberto; e, entre as 164 unidades com dado suficiente para o teste, a taxa de
conversão não cai
sistematicamente com a depleção do estoque ($p = 0{,}48$), nem há saturação
(termo quadrático nulo, $p = 0{,}92$). Quando a taxa é aproximadamente
constante, o fluxo passa a ser ditado pelo tamanho do estoque --- e é esse
resultado, e não a decomposição, que põe a fronteira onde ainda há Cerrado.
A Figura~\ref{fig:fronteira} decompõe a queda entre estoque e taxa.

\begin{figure}[htb]
\centering
\caption[Decomposição da mudança do fluxo entre os Atos II e III]{Decomposição da mudança do fluxo de conversão entre os Atos~II e
III, por região (a), e o estoque convertível remanescente de cada região como
percentual do seu próprio estoque de 1985 (b).}
\label{fig:fronteira}
\includegraphics[width=\textwidth]{cap4_fronteira_oferta.pdf}
\fontefig{elaboração própria (Pipeline~\#39). A parcela residual reúne tudo
o que não é o tamanho do estoque --- propensão a converter, atrito de acesso,
proteção e a troca de fonte de terra ---, e \textbf{não é uma medida de
demanda}: lê-la como tal inverteria o sentido do argumento, que se monta por
eliminação e por regularidade, e não pelo tamanho da parcela de estoque.}
\end{figure}

\subsection{O estoque ainda exposto à conversão}

Comparados em hectares, os estoques escondem o que interessa, porque o Sul
sempre foi a menor das três regiões. Postos na mesma régua --- cada região
como percentual do que tinha em 1985 --- as trajetórias se separam. O Sul
desce mais rápido e mais cedo, é o único que cruza os 60\%, e estabiliza a
partir de 2019: perde dois pontos em cinco anos, contra 3,5 do Centro e 3,4
do Norte, que seguem descendo juntos. Não é que o Sul tenha passado a
preservar: é que restou pouco Cerrado em pé a ser suprimido.

Em hectares, restam 6,56~Mha de vegetação natural ainda exposta à conversão
em Goiás, cerca de 60\% do que havia em 1985, e o Norte concentra 44\%
deles. O termo ``convertível'' mede exposição, e não disponibilidade: é o
Cerrado que ainda pode ser perdido, não uma reserva de terra que este
trabalho trate como aproveitável.

A definição adotada --- formação savânica e campo nativo --- não é
arbitrária, e decorre do que a própria depleção mostra: entre 1985 e 2024, a
formação savânica perdeu 38,2\% de sua área e o campo nativo 42,5\%,
enquanto a floresta perdeu 22,8\%. Savana e campo são as classes que a
fronteira de fato consome. O limite da medida é declarado: sem cadastro
ambiental integrado pixel a pixel, o estoque inclui Reserva Legal e Área de
Preservação Permanente (APP) que
na prática não podem ser convertidas, de modo que ele é um \emph{teto} --- a
terra realmente exposta é menor que 6,56~Mha, nunca maior. Um cenário de
sensibilidade descontando 20\% de Reserva Legal foi executado e não altera o
ordenamento entre regiões.

Esse limite merece ser levado às últimas consequências, porque nele mora a
alternativa mais séria à leitura de teto de oferta. O que este trabalho
elimina com dado é a \emph{proteção integral}: unidades de conservação de
uso restrito cobrem parcela pequena do estado, estavam congeladas desde
antes da freada e não explicam o momento em que o Sul desacelerou. Reserva
Legal e APP são outra coisa. Elas incidem dentro de cada propriedade,
ganharam fiscalização efetiva com o Cadastro Ambiental Rural a partir de
2012 --- justamente na janela de interesse --- e tendem a morder mais forte
onde o remanescente já é escasso, porque ali o que sobrou de vegetação
tende a ser precisamente o passivo legal. No Sul depletado, portanto,
``acabou a terra convertível'' e ``o que restou é legalmente inconversível''
produzem a mesma série observada. O cenário de $-20\%$ testado acima
compara regiões entre si e não separa essas duas leituras dentro do Sul;
separá-las exigiria o cadastro ambiental integrado pixel a pixel, que este
trabalho não teve. Fica registrado como limite, e não como questão vencida.

Uma fronteira pode parar por duas razões muito diferentes --- porque a terra
acabou ou porque a lei a barrou --- e as duas produzem o mesmo gráfico
enquanto pedem políticas opostas. Em Goiás, a proteção não estava no caminho
da fronteira. Se o teto fosse institucional, a proteção teria vindo antes da
freada e estaria onde a fronteira vai; não veio, e não está. Dos 6,56~Mha
ainda expostos, 6,35 estão fora de Unidade de Conservação de Proteção
Integral: 97\% pela malha de unidades de conservação e 94,3\% quando a mesma
conta é refeita pixel a pixel, que é a régua mais exigente. A Proteção
Integral cobre menos de 3\% em qualquer região, e congelou: 76\% da que Goiás
tem hoje já existia em 1985 e 89\% já existia em 2000 --- em quarenta anos
ela cresceu 0,12~Mha, contra 4,11~Mha suprimidos do mesmo estoque convertível
(os 10,67~Mha de savana e campo de 1985 reduzidos aos 6,56~Mha de hoje). O
que cresceu nos anos 2000 foi o Uso Sustentável, que admite uso rural e não
veda conversão.

O sentido de ``desprotegido'' aqui é estreito e precisa ser dito: fora da
única categoria que veda a conversão. Não significa livre para desmatar ---
Reserva Legal e APP seguem valendo dentro de cada propriedade e protegem
parte desse mesmo estoque, e é por isso que os 6,56~Mha são um teto. A conta
mede outra coisa: se o instrumento de política que poderia deter a fronteira
de propósito está no caminho dela. Não está --- e o Norte, que guarda 44\%
do que resta, não tem pela frente área protegida que barre a conversão.

\subsection{Custo de carbono e desenvolvimento municipal}

Duas contas fecham a frente, e nenhuma é econômica no sentido usual.

A primeira é de carbono. A conta segue a lógica de diferença de estoque do
IPCC em seu nível mais simples (\emph{Tier}~1): a cada hectare convertido
atribui-se a densidade de carbono da formação que havia ali, e a emissão é a
diferença entre o estoque antes e depois. Duas declarações de escopo são
devidas. A primeira é que a conta cobre apenas a \textbf{biomassa} --- acima
e abaixo do solo ---, e \emph{não} o carbono do solo, que no Cerrado é
parcela grande do estoque total; o número abaixo é, portanto, um piso do
passivo, não o passivo. A segunda é que a faixa reportada não é intervalo de
confiança estatístico: ela vem de rodar a mesma conta sob três conjuntos de
densidades por formação --- baixo, central e alto (decisão D18) ---,
colhidos da literatura de estoques do Cerrado
\cite{Bustamante2012, Grace2006}, das diretrizes do \citeonline{IPCC2006} e,
para a delimitação das fitofisionomias a que cada densidade se aplica, de
\citeonline{RibeiroWalter2008}. Vale registrar que o Cerrado tem razão raiz:parte
aérea alta, de modo que savana e campo estão longe de ser desprezíveis na
conta. Assim precificada, a conversão comprometeu da ordem de
973 milhões de toneladas (Mt) de dióxido de carbono equivalente
(CO\textsubscript{2}e), em faixa de 751 a 1.208 conforme o cenário de
densidade. Duas coisas surpreendem. A primeira é
quem paga: a floresta perde 2,6 vezes menos área que a formação savânica e
ainda assim emite mais (499 contra 458~Mt), porque é cerca de três vezes
mais densa --- o custo de carbono não é proporcional à área. A segunda é
quando se pagou: 80\% do total, 774~Mt, saiu no Ato~I, a 51,6~Mt por ano,
com o Sul como maior emissor do período --- 307~Mt, ou 40\% do ato, a
20,5~Mt por ano. Depois de 2001 o ritmo cai cerca de seis vezes (8,2~Mt por
ano no Ato~II) e a geografia se inverte --- no Ato~III o Sul é o menor
emissor (1,4~Mt/ano) e o Centro o maior (4,6). O centroide da emissão marcha
98~km ao norte, atrás da fronteira. No acumulado dos quarenta anos, note-se,
as três regiões emitem parcelas próximas (Sul 36\%, Norte 32\%, Centro
32\%): o que distingue o Sul não é o total, é o fato de ele ter emitido cedo.

A segunda conta é humana, e exige cuidado para não dizer mais do que mede.
Entre 2013 e 2021, a fronteira norte quase dobrou a área cultivada
($+93\%$, contra $+14\%$ no Sul). O desenvolvimento municipal, medido pelo
Índice FIRJAN de Desenvolvimento Municipal (IFDM) --- indicador composto de
emprego e renda, educação e saúde, publicado pela Federação das Indústrias do
Estado do Rio de Janeiro ---, subiu em toda parte no período --- em Goiás, de
0,49 a 0,63 --- mas
subiu no Norte na mesma medida que no Sul, e o Norte terminou onde começou
em relação a ele: o vão permanece intacto ($-0{,}083$ em 2023, IC 95\%
$[-0{,}108; -0{,}058]$, contra $-0{,}092$ em 2013). Quase dobrar a área não
comprou desenvolvimento adicional nem aproximou a fronteira do núcleo. Num
painel de efeitos fixos dentro do município, dobrar a área vale 0,008 ponto
de IFDM, contra os cerca de 0,15 que o índice subiu na década. A leitura é
associativa, com o gradiente espacial controlado, e não causal.

Vários dos resultados centrais deste capítulo são nulos, e um nulo mal
enunciado converte-se numa afirmação maior do que a evidência sustenta; o
alcance e os limites de cada um estão consolidados na
Seção~\ref{sec:disc-limites}. Os resultados que teriam sido publicados se um
teste específico não tivesse sido feito, e a distância entre cada versão
superada e a que ficou no lugar, estão no Apêndice~\ref{ap:superados}.
